% GENERATED FILE. Edit canonical lessons, then run npm run generate:books.
% canonical-source-hash: fnv1a64:fd074dd18b6b9766
% canonical-lessons: JA-C334-fukushi, JA-C334-kifu-suru, JA-C334-bokin, JA-C334-sewa, JA-C334-atsukau, JA-R334-review-329-334

\chapter{Helping the Community}
\label{ch:ja-helping-the-community}
\hlchaptermodality{334}

\begin{chapteropening}
\textbf{By the end of this chapter:} \emph{I can say how people help their community, and why.}

Complete the last lesson of chapter 334: I can say how people help their community, and why.
\end{chapteropening}

\section[\texorpdfstring{\ja{ふくし} (fukushi)}{ふくし (fukushi)}]{\ja{ふくし} (fukushi) — welfare; social services}

\label{lesson:JA-C334-fukushi}



\begin{quote}
Before the new one: say the Japanese for to give care, then the Japanese for bringing up children.
\end{quote}

\subsection*{You'll want to know: \ja{ふくし}}
\textbf{\ja{ふくし}} — \emph{fukushi} — \textquotedblleft{}welfare; social services\textquotedblright{}.

\begin{quote}
\textbf{\ja{こうれいしゃが} \ja{ふえたので、ふくしが} \ja{ひつようです。}} — \emph{kōreisha ga fueta node, fukushi ga hitsuyō desu} — \textquotedblleft{}The elderly have grown in number, so welfare is needed.\textquotedblright{}
\end{quote}

\begin{cousinweb}
\textbf{\ja{ふくし}} (\emph{fukushi}) is Sino-Japanese (kango): two Chinese characters read with their on'yomi, readings borrowed from Chinese. They are \emph{fuku}, \textquotedblleft{}blessing\textquotedblright{} and \emph{shi}, \textquotedblleft{}welfare\textquotedblright{} (Middle Chinese \emph{trhiX}) (KANJIDIC2; Middle Chinese in Baxter's transcription).
\end{cousinweb}

\begin{grammarlens}[title={What you've built}]
One more word for giving reasons and opinions.
\end{grammarlens}

\subsection*{Your turn}
\begin{itemize}
\raggedright
  \item \emph{Say it:} \emph{fukushi}
  \item \emph{Say it:} \emph{fukushi}, once more
  \item \emph{Say it:} \emph{kōreisha ga fueta node, fukushi ga hitsuyō desu} — The elderly have grown in number, so welfare is needed.
  \item \emph{Recall:} say the Japanese for to give care, then the Japanese for bringing up children, then say \emph{fukushi} again
\end{itemize}

\begin{tcolorbox}[breakable,colback=teal!4,colframe=teal!35!black,title={Before you move on}]
What does \ja{ふくし} mean? (\textquotedblleft{}Welfare; social services\textquotedblright{}.) Say it once more.
\end{tcolorbox}
\section[\texorpdfstring{\ja{きふする} (kifu suru)}{きふする (kifu suru)}]{\ja{きふする} (kifu suru) — to donate}

\label{lesson:JA-C334-kifu-suru}



\begin{quote}
Before the new one: say the Japanese for bringing up children, then the Japanese for welfare.
\end{quote}

\subsection*{You'll want to know: \ja{きふする}}
\textbf{\ja{きふする}} — \emph{kifu suru} — \textquotedblleft{}to donate\textquotedblright{}.

\begin{quote}
\textbf{\ja{たいふうで} \ja{みんなが} \ja{こまったので、おかねを} \ja{きふしました。}} — \emph{taifū de minna ga komatta node, okane o kifu shimashita} — \textquotedblleft{}The typhoon left everyone in trouble, so I gave money.\textquotedblright{}
\end{quote}

Polite \textbf{\ja{きふします}}, \emph{kifushimasu}; a \ja{する} verb. Past (polite): \textbf{\ja{ちちが} \ja{おかねを} \ja{きふしました。}}, \emph{chichi ga okane o kifu shimashita}.

\begin{cousinweb}
\textbf{\ja{きふする}} (\emph{kifu suru}) is built from a Sino-Japanese noun and the native verb \emph{suru}, \textquotedblleft{}to do\textquotedblright{}. The noun, \emph{kifu}, is two Chinese characters read with their on'yomi, readings borrowed from Chinese: \emph{ki}, \textquotedblleft{}draw near\textquotedblright{} (Middle Chinese \emph{kjeH}) and \emph{fu}, \textquotedblleft{}attach\textquotedblright{} (\emph{pjuH}) (KANJIDIC2; JMdict; Middle Chinese in Baxter's transcription).
\end{cousinweb}

\begin{grammarlens}[title={What you've built}]
One more word for giving reasons and opinions.
\end{grammarlens}

\subsection*{Your turn}
\begin{itemize}
\raggedright
  \item \emph{Say it:} \emph{kifu suru}
  \item \emph{Say it:} \emph{kifu suru}, once more
  \item \emph{Say it:} \emph{taifū de minna ga komatta node, okane o kifu shimashita} — The typhoon left everyone in trouble, so I gave money.
  \item \emph{Recall:} say the Japanese for bringing up children, then the Japanese for welfare, then say \emph{kifu suru} again
\end{itemize}

\begin{tcolorbox}[breakable,colback=teal!4,colframe=teal!35!black,title={Before you move on}]
What does \ja{きふする} mean? (\textquotedblleft{}To donate\textquotedblright{}.) Say it once more.
\end{tcolorbox}
\section[\texorpdfstring{\ja{ぼきん} (bokin)}{ぼきん (bokin)}]{\ja{ぼきん} (bokin) — fund-raising; a collection}

\label{lesson:JA-C334-bokin}



\begin{quote}
Before the new one: say the Japanese for welfare, then the Japanese for to donate.
\end{quote}

\subsection*{You'll want to know: \ja{ぼきん}}
\textbf{\ja{ぼきん}} — \emph{bokin} — \textquotedblleft{}fund-raising; a collection\textquotedblright{}.

\begin{quote}
\textbf{\ja{たいふうが} \ja{あったので、えきで} \ja{ぼきんを} \ja{しています。}} — \emph{taifū ga atta node, eki de bokin o shite imasu} — \textquotedblleft{}There was a typhoon, so people are collecting money at the station.\textquotedblright{}
\end{quote}

\begin{cousinweb}
\textbf{\ja{ぼきん}} (\emph{bokin}) is Sino-Japanese (kango): two Chinese characters read with their on'yomi, readings borrowed from Chinese. They are \emph{bo}, \textquotedblleft{}recruit\textquotedblright{} (Middle Chinese \emph{muH}) and \emph{kin}, \textquotedblleft{}gold\textquotedblright{} (\emph{kim}) (KANJIDIC2; Middle Chinese in Baxter's transcription).
\end{cousinweb}

\begin{grammarlens}[title={What you've built}]
One more word for giving reasons and opinions.
\end{grammarlens}

\subsection*{Your turn}
\begin{itemize}
\raggedright
  \item \emph{Say it:} \emph{bokin}
  \item \emph{Say it:} \emph{bokin}, once more
  \item \emph{Say it:} \emph{taifū ga atta node, eki de bokin o shite imasu} — There was a typhoon, so people are collecting money at the station.
  \item \emph{Recall:} say the Japanese for welfare, then the Japanese for to donate, then say \emph{bokin} again
\end{itemize}

\begin{tcolorbox}[breakable,colback=teal!4,colframe=teal!35!black,title={Before you move on}]
What does \ja{ぼきん} mean? (\textquotedblleft{}Fund-raising; a collection\textquotedblright{}.) Say it once more.
\end{tcolorbox}
\section[\texorpdfstring{\ja{せわ} (sewa)}{せわ (sewa)}]{\ja{せわ} (sewa) — looking after; care}

\label{lesson:JA-C334-sewa}



\begin{quote}
Before the new one: say the Japanese for to donate, then the Japanese for fund-raising.
\end{quote}

\subsection*{You'll want to know: \ja{せわ}}
\textbf{\ja{せわ}} — \emph{sewa} — \textquotedblleft{}looking after; care\textquotedblright{}.

\begin{quote}
\textbf{\ja{いぬの} \ja{せわを} \ja{するから、はやく} \ja{かえります。}} — \emph{inu no sewa o suru kara, hayaku kaerimasu} — \textquotedblleft{}I look after the dog, so I go home early.\textquotedblright{}
\end{quote}

\begin{cousinweb}
\textbf{\ja{せわ}} (\emph{sewa}) is Sino-Japanese (kango): two Chinese characters read with their on'yomi, readings borrowed from Chinese. They are \emph{se}, \textquotedblleft{}world\textquotedblright{} (Middle Chinese \emph{syejH}) and \emph{wa}, \textquotedblleft{}talk\textquotedblright{} (\emph{hwaejH}) (KANJIDIC2; Middle Chinese in Baxter's transcription).
\end{cousinweb}

\begin{grammarlens}[title={What you've built}]
One more word for giving reasons and opinions.
\end{grammarlens}

\subsection*{Your turn}
\begin{itemize}
\raggedright
  \item \emph{Say it:} \emph{sewa}
  \item \emph{Say it:} \emph{sewa}, once more
  \item \emph{Say it:} \emph{inu no sewa o suru kara, hayaku kaerimasu} — I look after the dog, so I go home early.
  \item \emph{Recall:} say the Japanese for to donate, then the Japanese for fund-raising, then say \emph{sewa} again
\end{itemize}

\begin{tcolorbox}[breakable,colback=teal!4,colframe=teal!35!black,title={Before you move on}]
What does \ja{せわ} mean? (\textquotedblleft{}Looking after; care\textquotedblright{}.) Say it once more.
\end{tcolorbox}
\section[\texorpdfstring{\ja{あつかう} (atsukau)}{あつかう (atsukau)}]{\ja{あつかう} (atsukau) — to handle; to treat}

\label{lesson:JA-C334-atsukau}



\begin{quote}
Before the new one: say the Japanese for fund-raising, then the Japanese for looking after.
\end{quote}

\subsection*{You'll want to know: \ja{あつかう}}
\textbf{\ja{あつかう}} — \emph{atsukau} — \textquotedblleft{}to handle; to treat\textquotedblright{}.

\begin{quote}
\textbf{\ja{この} \ja{おさらは} \ja{たかいから、こどもは} \ja{あつかわないで} \ja{ください。}} — \emph{kono osara wa takai kara, kodomo wa atsukawanaide kudasai} — \textquotedblleft{}This plate is expensive, so children, please do not handle it.\textquotedblright{}
\end{quote}

Polite \textbf{\ja{あつかいます}}, \emph{atsukaimasu}; a godan verb. Past (polite): \textbf{\ja{てんいんが} \ja{にもつを} \ja{あつかいました。}}, \emph{ten'in ga nimotsu o atsukaimashita}.

\begin{cousinweb}
\textbf{\ja{あつかう}} (\emph{atsukau}) is a native Japanese word (wago), inherited rather than borrowed from Chinese, and its older origin is not settled. Its kanji was attached to it later as a kun'yomi, a native reading given to a Chinese character (JMdict, KANJIDIC2).
\end{cousinweb}

\begin{grammarlens}[title={What you've built}]
One more word for giving reasons and opinions.
\end{grammarlens}

\subsection*{Your turn}
\begin{itemize}
\raggedright
  \item \emph{Say it:} \emph{atsukau}
  \item \emph{Say it:} \emph{atsukau}, once more
  \item \emph{Say it:} \emph{kono osara wa takai kara, kodomo wa atsukawanaide kudasai} — This plate is expensive, so children, please do not handle it.
  \item \emph{Recall:} say the Japanese for fund-raising, then the Japanese for looking after, then say \emph{atsukau} again
\end{itemize}

\begin{tcolorbox}[breakable,colback=teal!4,colframe=teal!35!black,title={Before you move on}]
What does \ja{あつかう} mean? (\textquotedblleft{}To handle; to treat\textquotedblright{}.) Say it once more.
\end{tcolorbox}
\section[(dialogue)]{Looking back: from Common Sense to Helping the Community}

\label{lesson:JA-R334-review-329-334}



\begin{quote}
Say the Japanese for common sense and only natural.
\end{quote}

\subsection*{Your turn}
Say these aloud:

\begin{itemize}
\raggedright
  \item common sense, only natural, naturally, one's turn, a cigarette end
  \item as one should, firmly, to follow an example, to line up neatly, to revise
  \item society, a citizen, a resident, politics, a law
  \item a right, a duty, equal, discrimination, all members
  \item young people, an elderly person, population, to give care, bringing up children
  \item welfare, to donate, fund-raising, looking after, to handle
\end{itemize}

\begin{tcolorbox}[breakable,colback=teal!4,colframe=teal!35!black,title={Before you move on}]
Common sense? (\textbf{\ja{じょうしき}}, \emph{jōshiki}.) Only natural? (\textbf{\ja{あたりまえ}}, \emph{atarimae}.) Naturally? (\textbf{\ja{とうぜん}}, \emph{tōzen}.) One's turn? (\textbf{\ja{じゅんばん}}, \emph{junban}.) A cigarette end? (\textbf{\ja{すいがら}}, \emph{suigara}.) As one should? (\textbf{\ja{ちゃんと}}, \emph{chanto}.) Firmly? (\textbf{\ja{しっかり}}, \emph{shikkari}.) To follow an example? (\textbf{\ja{みならう}}, \emph{minarau}.) To line up neatly? (\textbf{\ja{そろえる}}, \emph{soroeru}.) To revise? (\textbf{\ja{あらためる}}, \emph{aratameru}.) Society? (\textbf{\ja{しゃかい}}, \emph{shakai}.) A citizen? (\textbf{\ja{しみん}}, \emph{shimin}.) A resident? (\textbf{\ja{じゅうみん}}, \emph{jūmin}.) Politics? (\textbf{\ja{せいじ}}, \emph{seiji}.) A law? (\textbf{\ja{ほうりつ}}, \emph{hōritsu}.) A right? (\textbf{\ja{けんり}}, \emph{kenri}.) A duty? (\textbf{\ja{ぎむ}}, \emph{gimu}.) Equal? (\textbf{\ja{びょうどう}}, \emph{byōdō}.) Discrimination? (\textbf{\ja{さべつ}}, \emph{sabetsu}.) All members? (\textbf{\ja{ぜんいん}}, \emph{zen'in}.) Young people? (\textbf{\ja{わかもの}}, \emph{wakamono}.) An elderly person? (\textbf{\ja{こうれいしゃ}}, \emph{kōreisha}.) Population? (\textbf{\ja{じんこう}}, \emph{jinkō}.) To give care? (\textbf{\ja{かいごする}}, \emph{kaigo suru}.) Bringing up children? (\textbf{\ja{こそだて}}, \emph{kosodate}.) Welfare? (\textbf{\ja{ふくし}}, \emph{fukushi}.) To donate? (\textbf{\ja{きふする}}, \emph{kifu suru}.) Fund-raising? (\textbf{\ja{ぼきん}}, \emph{bokin}.) Looking after? (\textbf{\ja{せわ}}, \emph{sewa}.) To handle? (\textbf{\ja{あつかう}}, \emph{atsukau}.)
\end{tcolorbox}
